\section{The predication morphism generators}

\begin{proposition}[Converse]\label{thm:predication-converse}\lean{Infomaterial.PredicationCell.converse}\leanok\uses{thm:predication}
Converse swaps subject and object while preserving the predicate. It is the
orientation-reversing generator on predication cells.
\end{proposition}

\begin{proposition}[Projections]\label{thm:predication-proj}\lean{Infomaterial.PredicationCell.subjProj}\leanok\uses{thm:predication}
The subject projection sends $(s,p,o)$ to $(s,p,s)$, and the object
projection sends it to $(o,p,o)$. Both collapse an oriented cell onto a
reflexive diagonal.
\end{proposition}

\begin{proposition}[Specialization]\label{thm:predication-specialize}\lean{Infomaterial.PredicationCell.specialize}\leanok\uses{thm:predication}
Specialization by a predicate move $f$ keeps the entity endpoints fixed and
acts on the middle coordinate: $(s,p,o)\mapsto (s,f(p),o)$.
\end{proposition}

\begin{theorem}[Converse is an involution]\label{thm:converse-involution}\lean{Infomaterial.PredicationCell.converse_converse}\leanok\uses{thm:predication-converse}
Applying converse twice returns the original predication cell. The converse
generator has order two.
\end{theorem}

\begin{theorem}[The seed is self-converse]\label{thm:converse-seed}\lean{Infomaterial.PredicationCell.converse_seed}\leanok\uses{thm:predication-converse, thm:predication-seed}
The seed of any utterance is fixed by converse, because its subject and object
are the same token.
\end{theorem}

\begin{theorem}[Orientation is real]\label{thm:converse-orientation}\lean{Infomaterial.PredicationCell.converse_ne}\leanok\uses{thm:predication-converse}
If a predication cell has distinct subject and object, then its converse is
not equal to it. Off the diagonal, orientation is information.
\end{theorem}

\begin{proposition}[Projections are idempotent]\label{thm:subjproj-idem}\lean{Infomaterial.PredicationCell.subjProj_idem}\leanok\uses{thm:predication-proj}
Subject and object projections are idempotent. Once a cell has been collapsed
to the corresponding diagonal, applying the same projection again changes
nothing.
\end{proposition}

\begin{theorem}[Specialization is functorial]\label{thm:specialize-comp}\lean{Infomaterial.PredicationCell.specialize_comp}\leanok\uses{thm:predication-specialize}
Specialization composes by composition of predicate moves, and identity
specialization is the identity on predication cells. The predicate coordinate
therefore carries the functorial action generated by predicating.
\end{theorem}
